\documentclass[11pt]{article}
\usepackage[margin=1in]{geometry}
\usepackage{booktabs}
\usepackage{xcolor}
\usepackage{listings}
\usepackage{enumitem}
\usepackage{titlesec}
\usepackage[hidelinks]{hyperref}

\definecolor{codebg}{gray}{0.96}
\lstset{
  basicstyle=\ttfamily\small,
  backgroundcolor=\color{codebg},
  breaklines=true,
  frame=single,
  framesep=4pt,
  columns=fullflexible,
  keepspaces=true,
  showstringspaces=false,
}
\titleformat{\section}{\large\bfseries}{\thesection}{0.6em}{}
\setlist{nosep,leftmargin=1.4em}

\title{\textbf{R730xd Backup Node: Build Plan}}
\author{NetFRAME Home Lab}
\date{Drafted 2026-07-22. READ AND FILL THE TBD FIELDS before executing.}

\begin{document}
\maketitle

\section{Purpose and role}
A Dell PowerEdge R730xd added as a second, independent storage node whose job is twofold:
\begin{enumerate}
  \item \textbf{Backup target.} It runs its own Proxmox Backup Server (PBS) and \emph{pulls}
  a copy of Randy's PBS datastore over a scheduled sync job. This gives a real backup of the
  backups and closes the audit finding that all backups currently sit only on Randy (one node,
  one rack, one UPS bus). See \texttt{audit/findings/30-gap-analysis.md} item R-04.
  \item \textbf{Transcode node.} It holds the Jellyfin video-transcode GPU, taking that role
  off Randy.
\end{enumerate}
The sync direction is pull: the backup node holds credentials to read Randy, and Randy stays
passive. That is the safer posture (a compromise of the backup node cannot write to Randy).

\section{Fields to confirm before you start (TBD)}
\begin{tabular}{@{}ll@{}}
\toprule
Field & Value \\
\midrule
Node hostname & \textbf{TBD} (pick a name in your scheme) \\
GPU card & \textbf{TBD}: RX 580 moved from Randy (VAAPI), or a new NVIDIA card (NVENC) \\
Mgmt IP (VLAN 1) & \texttt{192.168.10.188} (confirm free in Kea + Pi-hole) \\
Service IP (VLAN 30, 10G) & \texttt{192.168.30.188} (confirm free) \\
iDRAC IP (VLAN 20) & \texttt{192.168.20.23} (next after Randy .22; confirm free) \\
Rack location & \textbf{TBD}: a different rack buys real fault isolation for the backup \\
Sync scope & Randy datastore \texttt{datastore} only (skip the degraded \texttt{bulk} media) \\
\bottomrule
\end{tabular}

\section{Hardware}
\begin{itemize}
  \item Chassis: R730xd, 24x 2.5" front bays (plus rear flexbay).
  \item Drives: 7x 1.8\,TB SAS, 15x 1.2\,TB SAS, 2x 250\,GB SSD (boot).
  \item Storage controller: \textbf{HBA330 Mini in IT mode} (or a PERC H730 flashed to IT mode),
  so ZFS owns redundancy. Do not use hardware RAID. This matches Randy's IT-mode approach.
  \item NIC: a 10\,GbE port for the VLAN 30 service path (ConnectX or X520 class), plus the
  onboard 1\,G for VLAN 1 management, mirroring Randy's dual-homed layout.
  \item GPU: single-slot or dual-slot card in a PCIe riser; confirm the R730xd riser and power
  support the chosen card (the RX 580 needs a 6-pin/8-pin PCIe power lead).
\end{itemize}

\section{Network (dual-homed, like Randy)}
\begin{itemize}
  \item Management, corosync (if clustered), and the web UIs ride VLAN 1 on the onboard 1\,G.
  \item PBS sync traffic and any NFS ride VLAN 30 on the 10\,G NIC, reaching Randy directly at
  \texttt{192.168.30.187} at 10\,G. Keep the PBS sync on the \texttt{.30} path so it never
  touches the management VLAN.
  \item iDRAC on VLAN 20 (tagged), reachable from Ares \texttt{enp0s31f6.20}, factory creds
  rotated into Vaultwarden (same policy as the other BMCs).
\end{itemize}

\section{OS and boot}
Install Proxmox VE onto the two 250\,GB SSDs as a \textbf{ZFS mirror} (\texttt{rpool}) using the
PVE installer (select the two SSDs, RAID1). PBS installs on top as a service package, the same
pattern Randy already runs.
\begin{lstlisting}[language=bash]
# after PVE is installed and updated, add the PBS server package:
apt update && apt install proxmox-backup-server
\end{lstlisting}

\section{ZFS data pool (Option B: 14 active + 1 hot spare)}
One pool with three RAIDZ2 vdevs plus a hot spare. Always use \texttt{/dev/disk/by-id} names,
never \texttt{sdX} (which reorder across reboots).
\begin{lstlisting}[language=bash]
# Replace the by-id placeholders with the real device ids (ls -l /dev/disk/by-id | grep -i sas)
zpool create -o ashift=12 \
  -O compression=zstd -O atime=off -O xattr=sa -O relatime=on \
  vault \
  raidz2 <7x 1.8TB by-id>            `# vdev A: ~8.2 TiB usable` \
  raidz2 <7x 1.2TB by-id>            `# vdev B1: ~5.4 TiB usable` \
  raidz2 <7x 1.2TB by-id>            `# vdev B2: ~5.4 TiB usable` \
  spare  <1x 1.2TB by-id>           `# hot spare for the 1.2TB vdevs`

# PBS datastore dataset, larger recordsize suits PBS chunk files:
zfs create -o recordsize=1M vault/pbs
zpool status vault
\end{lstlisting}
Usable capacity is about 19\,TiB, enough to hold Randy's \texttt{datastore} PBS data with room
to grow. Rename the pool \texttt{vault} to your preference if you like.

\subsection*{Spare coverage caveat (important)}
The single 1.2\,TB hot spare can auto-replace any of the fourteen 1.2\,TB drives, but it
\textbf{cannot} replace a 1.8\,TB drive (a spare must be at least as large as the disk it
replaces). Vdev A (the 1.8\,TB drives) therefore has no automatic spare. Keep a cold 1.8\,TB
drive on the shelf and replace vdev-A failures by hand.

\section{PBS datastore}
\begin{lstlisting}[language=bash]
proxmox-backup-manager datastore create backup /vault/pbs
# local retention (independent of Randy) + verify + gc:
proxmox-backup-manager prune-job create backup-prune \
  --store backup --schedule 'daily' \
  --keep-daily 14 --keep-weekly 8 --keep-monthly 6
proxmox-backup-manager verify-job create backup-verify \
  --store backup --schedule 'sun 07:00'
proxmox-backup-manager garbage-collection ... # schedule daily GC in the UI or datastore config
\end{lstlisting}

\section{PBS remote sync (pull from Randy)}
\subsection*{On Randy (source): a read-only sync identity}
\begin{lstlisting}[language=bash]
proxmox-backup-manager user create sync@pbs
proxmox-backup-manager user generate-token sync@pbs backup-node
#   -> copy the returned token secret into Vaultwarden
proxmox-backup-manager acl update /datastore/datastore DatastoreReader \
  --auth-id 'sync@pbs!backup-node'
# note Randy's fingerprint for the remote definition:
proxmox-backup-manager cert info | grep -i fingerprint
\end{lstlisting}

\subsection*{On the R730xd (target): the remote and the sync job}
\begin{lstlisting}[language=bash]
proxmox-backup-manager remote create randy \
  --host 192.168.30.187 \
  --auth-id 'sync@pbs!backup-node' \
  --password '<token-secret-from-vaultwarden>' \
  --fingerprint '<randy-fingerprint>'

# pull Randy's datastore into the local 'backup' datastore, nightly AFTER Randy's jobs finish.
# remove-vanished false: the backup keeps snapshots even after Randy prunes them.
proxmox-backup-manager sync-job create randy-pull \
  --remote randy --remote-store datastore \
  --store backup \
  --schedule '06:00' \
  --remove-vanished false
proxmox-backup-manager sync-job run randy-pull   # first run, verify it pulls
\end{lstlisting}
Randy's nightly guest backups run roughly 02:00 to 08:00 UTC; set the sync schedule to land
after they finish (adjust the \texttt{06:00} above to your local time as needed).

\section{GPU passthrough for Jellyfin}
Run Jellyfin in an unprivileged LXC with the GPU passed through. The exact steps depend on the
card, which is still TBD.
\begin{itemize}
  \item \textbf{RX 580 (AMD, VAAPI):} bind \texttt{/dev/dri} into the container. Jellyfin uses
  VAAPI for AMD transcode (this is simpler than ROCm, which is not needed for transcode). Ensure
  the amdgpu kernel module loads on the host.
  \item \textbf{New NVIDIA (NVENC):} install the pinned NVIDIA driver on the host plus the
  container toolkit, then expose the nvidia devices to the LXC. This mirrors the QuarkyLab and
  Jarvis GPU stack pattern (kernel pin plus held driver packages).
\end{itemize}
Jellyfin media can mount from Randy or from a dataset on this node; decide during setup.

\section{Verification checklist}
\begin{itemize}
  \item \texttt{zpool status vault} shows all vdevs ONLINE, the spare AVAIL, no errors.
  \item Reboot once, confirm the pool auto-imports and PVE boots from the SSD mirror.
  \item \texttt{ip -br a} shows the VLAN 30 10\,G leg at \texttt{.30.188}, and \texttt{ping
  192.168.30.187} reaches Randy at 10\,G.
  \item \texttt{proxmox-backup-manager sync-job run randy-pull} completes and snapshots appear in
  the local \texttt{backup} datastore.
  \item A manual \texttt{verify-job} on the local datastore passes.
  \item Jellyfin transcodes a file using the GPU (check \texttt{nvidia-smi} or \texttt{radeontop}
  during a transcode).
  \item Add the node's node-exporter to Prometheus and confirm the target is UP.
\end{itemize}

\section{Rollback and safety}
\begin{itemize}
  \item This node is purely additive. Nothing on Randy changes except the creation of a
  read-only \texttt{sync@pbs} token, which is revocable: \texttt{proxmox-backup-manager user
  delete 'sync@pbs'}.
  \item If the sync misbehaves, disable the job (\texttt{sync-job update randy-pull --schedule
  ''}) with no effect on Randy.
  \item The pool and datastore live only on this node; destroying and rebuilding them touches
  nothing else.
\end{itemize}

\section{Open follow-ons}
\begin{itemize}
  \item True offsite copy (Backblaze B2 or a remote PBS) as a third layer, later.
  \item Add this node to the monitoring coverage matrix and alert rules.
  \item Fill the TBD fields at the top before executing.
\end{itemize}

\end{document}
